\section{An application to polynomial moments}

A normalizer of the acting CM field can strengthen the rank-two
bound. Suppose a group of invertible geometric quasi-endomorphisms normalizes
\(E\). After defining representatives of its finite image in
\(\operatorname{Aut}(E)\) over a finite base extension, the
embeddings in the same orbit have isomorphic Frobenius components
at good split primes outside finitely many exceptions.
Consequently the nonordinary balanced pairs form a union of orbits
on conjugate pairs. For \(b=3\), for example, a transitive action
on its three pairs forces ordinarity. More generally, simultaneous
vanishing of coherent Frobenius in \(\lceil b/2\rceil\) distinct
rank-two balanced pairs contradicts Corollary~\ref{cor:ranktwo}.
On either side of a balanced pair the coherent space has
dimension one, so its zero Frobenius scalar makes the pair
nonordinary. A single isolated moment coefficient usually supplies
just one such scalar.

The arbitrary-rank result also gives a moment obstruction without
a symmetry of the cover. Write \(f\) on the sphere quadric
\(t^2+2xy=1\) as a Laurent polynomial in \(x,t\), by substituting
\(y=(1-t^2)/(2x)\). Its sphere moment is
\[
 \mathcal A(f^m)=\operatorname{CT}_x\,\frac12
                       \int_{-1}^1 f(x,t)^m\,dt ,
\]
interpreted as a rational linear functional on the Laurent chart.
Assign weight \(2i+j\) to \(x^it^j\).

\begin{proposition}[A squarefree sextic face]\label{prop:sextic}
Suppose the greatest Laurent weight of \(f\) is four, with
initial form \(t^4z^{-2}P(z)\), where \(z=x/t^2\) and \(P\)
has degree four, is squarefree, and has nonzero constant term.
Then over any characteristic-zero field some positive power
has nonzero sphere moment. All smaller Laurent weights are
allowed. If the phase has normal-form degree at most six,
its least sphere degree is exactly six.
\end{proposition}

\begin{proof}
If the moments all vanish, specialize the finitely generated
coefficient algebra over \(\Q\) to a number field after inverting
the leading and constant coefficients of \(P\) and its discriminant.
The vanishing identities and the specified initial form survive.
At a good prime of degree one with rational characteristic
\(p=4k+1\), the largest possible \(t\)-degree in the angular
constant term of \(f^k\) is \(4k=p-1\).
Integration of an even \(t\)-degree \(j\) divides by \(j+1\)
and odd degrees integrate to zero. Multiplying by \(p\) and
reducing modulo the prime therefore gives
\[
       p\mathcal A(f^k)\equiv [z^{2k}]P(z)^k\pmod p .
\]
Only the initial form can contribute to this top degree.

The normalized cyclic cover \(W^4=z^2P(z)\) has branch
exponents \((2,1,1,1,1,2)\), including infinity.
Its primitive order-four Jacobian factor has \(\Q(i)\)-rank
four and coherent signature \((1,3)\), by the fractional-part
formula. It is ordinary at the degree-one primes supplied by
Proposition~\ref{prop:signatures} after finite extension.
Here the deck action again supplies the polarization-compatible
field, and the degree-one split primes satisfy \(p=4k+1\).
The coherent \(H^1\) comparison of the curve with its Jacobian
is equivariant for the deck action.
For the character of \(W\), the coherent eigensheaf has generators
\(W\) on the finite chart and \(z^{-2}W\) at infinity. The
\v{C}ech quotient has basis the class of \(z^{-1}W\).
Its Frobenius image is
\(z^{-p}W^p=z^{-p+2k}P(z)^kW\); the coefficient of the basis class
is exactly \([z^{2k}]P(z)^k\). Ordinary Frobenius is bijective
on coherent \(H^1\), so this scalar is nonzero at these primes.
This contradicts purity.

In the normal form, the negative angular exponents come from
monomials involving \(y\), with at most one of \(x,y\) present.
The \(x^{-2}t^8\) term in the initial form comes from \(y^2t^4\)
with nonzero coefficient. This term has degree six and cannot
be supplied by a term of lower degree. Thus the stated normal-form
degree bound gives least sphere degree six.
\end{proof}

\subsection{From Hasse--Witt restraints to endpoint bounds}

At a split prime the stable rank of coherent Frobenius on an
embedding component is its slope-zero multiplicity. In
particular, the rank of the first Frobenius map bounds this
multiplicity from above \cite[Definition~2.1 and Lemma~2.2]{bouw}.
Using the complementary component gives an upper bound on the
unit-slope multiplicity as well.

\begin{proposition}[Testing simultaneous rank restraints]\label{prop:rankcaps}
Suppose a split mixed pair has endpoint restraints \(z\le R\)
and \(o\le R'\), obtained from coherent ranks on the corresponding
components. If the ordinary endpoints fit both restraints, use
them. Otherwise set
\[
       z=\min(R,h-m-1),\qquad o=\min(R',m-1).
\]
For all pairs, the simultaneous inequalities are numerically
compatible with the given endpoint restraints exactly when
these choices pass \eqref{eq:interval} with the strict unit
margin. If a definite pair violates its endpoint restraints
there is no compatible polygon.
\end{proposition}

\begin{proof}
The ordinary endpoints give zero excess. Once they are
unavailable, a nonordinary choice is necessary, and the proposed
values are coordinatewise largest among allowable endpoints.
Each interval in \eqref{eq:interval} can only shrink as either
endpoint multiplicity increases. Thus if the displayed choices
fail the test, every collection of slopes fails it. If they pass,
Theorem~\ref{thm:extremal} realizes the simultaneous numerical
inequalities.
\end{proof}

For \(h\ge3\), a vanishing coherent line on a pair of type
\((1,h-1)\) can be matched numerically by slopes
\(0^{h-2},1/2,1/2\) on its dimension-one side. The paired excess
is two at no index, and the ordinary slack of this pair is one at
every index. Thus if another mixed pair stays ordinary, the
chosen pair can fail while all inequalities hold. A test of one
line typically needs another computation. The next examples
supply one explicitly.

\subsection{Arbitrarily large supporting faces}

The prime-isolation argument also works with a prime power.
If the highest angular-constant height at a chosen moment order
is \(p^r-1\), and the Laurent coefficients are integral at the
prime, multiplying the moment by \(p^r\) isolates this highest
coefficient modulo the prime. Indeed every smaller height has
interval denominator of valuation less than \(r\).
For the Kummer model \(W^n=F(z)\) whose selected coherent line
has generator \(W/z\), Frobenius to the \(r\)-th power
acts on this line when \(p^r-1=nk\). Its scalar is
\([z^{p^r-1}]F(z)^k\), by projecting
\(z^{-p^r}W^{p^r}=z^{-p^r}F(z)^k W\) onto the generator.

\begin{proposition}[Powering a sextic face]\label{prop:scaled}
For every \(L\ge1\), the highest-weight form
\[
       t^{4L} z^{-2L}P(z)^L,\qquad z=x/t^2 ,
\]
with \(P\) as in Proposition~\ref{prop:sextic}, forces a nonzero
power moment, with arbitrary strictly lower-weight terms.
For phases of normal-form degree at most \(6L\) the least sphere
degree is exactly \(6L\).
\end{proposition}

\begin{proof}
Use the same connected conductor-four cover as in
Proposition~\ref{prop:sextic}; its Hasse coefficient is nonzero at
the good degree-one primes after one extension. Further adjoining
\(\zeta_{4L}\) keeps the conclusion at all but finitely many such
primes. With \(p=4Lk+1\), the moment at order \(k\) isolates
\([z^{2Lk}]P^{Lk}\), the original scalar at conductor-four
exponent \(Lk\). The negative endpoint has leading monomial
\(y^{2L}t^{4L}\) up to a nonzero scalar. Specialization from an
arbitrary characteristic-zero coefficient field works as in
Proposition~\ref{prop:sextic}.
\end{proof}

\begin{example}[A selected line on a nonordinary split cover]
Take \(P=(z-1)^2(z-5)^3\) and \(W^9=z^5P(z)\) over
\(\mathbb F_{73}\). Its primitive relative rank is two, with
three balanced pairs. The scalars on characters \(2\) and \(4\)
are zero and \(22\), respectively, by coefficient extraction
with the branch exponents reduced modulo nine. Explicitly, for
character \(i\) let \((a_0,a_1,a_5)\) be the least residues of
\((5i,2i,3i)\) modulo nine. The scalar is
\[
       [z^{72-8a_0}](z-1)^{8a_1}(z-5)^{8a_5}\pmod {73}.
\]
Thus the pair
\((2,7)\) fails and the pair \((4,5)\) is ordinary. Its unit
trace ensures that the relative trace norm is nonzero.
Two bad pairs would give norm divisibility by \(73^4\),
incompatible with the Weil bound \(64\cdot73^3\). This makes
the pair \((1,8)\) ordinary, so the selected scalar on character
one is nonzero, while the primitive factor is nonordinary.

Let \(73\nmid L\), and take highest weight \(9L\) with form
\(t^{9L}z^{-4L}P(z)^L\). With coefficients integral at 73,
all strictly lower weights are allowed, and the top form need
only have the specified reduction. Choose \(r\) with
\(73^r\equiv1\pmod {9L}\), and take moment order
\((73^r-1)/(9L)\). Prime-power isolation extracts the
\(r\)-fold scalar on the original connected cover.
It is nonzero, so the moment is nonzero. With normal-form
degree at most \(13L\), the negative endpoint gives
least sphere degree \(13L\). The exact coefficient
replay also gives the selected scalar \(11\bmod73\).
\end{example}

\begin{example}[Increasing connected conductor]
Set \(N=(37^r-1)/4\) for \(r\ge2\) and
\(P=(z-1)^2(z-3)^3\). At highest weight \(N\) with initial
form \(t^N z^{-4}P(z)\), the fourth moment is nonzero for all
strictly lower terms integral at 37. In fact it gives
\([z^{16}]P^4\equiv32\pmod {37}\) after multiplication by
\(37^r\). The connected cover \(W^N=z^{N-4}P(z)\) has this
selected scalar for Frobenius to the \(r\)-th power.

The primitive factor is nonordinary at 37. For a primitive
character \(i\), put \(x=i/N\). Its coherent dimension is
\(1+\lfloor4x\rfloor-\lfloor2x\rfloor-\lfloor3x\rfloor\);
this is two for \(1/4<x<1/3\) and one for \(0<x<1/4\).
Since \(9\mid N\), the character \(i=N/3-1\) is primitive.
It has ratio in \((1/4,1/3)\), while the representative
\(37^{r-1}i\equiv N/3-37^{r-1}\pmod N\) has ratio in
\((0,1/4)\). The unequal dimensions in a Frobenius orbit
prevent coherent Frobenius on the factor from being invertible.
The selected line still has nonzero iterated scalar. The
negative endpoint gives least sphere degree \(N+4\) if the
normal-form degree is at most \(N+4\).
\end{example}
